\documentclass[a4paper,11pt]{amsart}

\usepackage[margin=1in]{geometry}
\usepackage[T1]{fontenc}
\usepackage{amsmath,amssymb,mathtools}
\usepackage[protrusion=true,expansion=false]{microtype}

\newtheorem{theorem}{Theorem}[section]
\newtheorem{proposition}[theorem]{Proposition}
\newtheorem{lemma}[theorem]{Lemma}
\theoremstyle{remark}
\newtheorem{remark}[theorem]{Remark}

\newcommand{\Z}{\mathbb Z}
\newcommand{\Q}{\mathbb Q}
\newcommand{\C}{\mathbb C}
\newcommand{\PP}{\mathbb P}
\DeclareMathOperator{\Comp}{Comp}
\DeclareMathOperator{\GL}{GL}
\DeclareMathOperator{\Sp}{Sp}
\DeclareMathOperator{\Lie}{Lie}
\DeclareMathOperator{\rank}{rank}
\DeclareMathOperator{\im}{im}
\DeclareMathOperator{\Span}{span}

\numberwithin{equation}{section}

\title{Even All-Half Thinness}
\date{}

\begin{document}

\begin{abstract}
We prove thinness for the even all-half hypergeometric tower
\[
        \Gamma_d=
        \left\langle
        \Comp((x-1)^d),\Comp((x+1)^d)
        \right\rangle,
        \qquad d\ge4\text{ even}.
\]
After a rational Levelt normal-form change, the group is generated by a regular unipotent symplectic matrix and a symplectic transvection.  The main input is a coefficient-cone renewal theorem for the odd kernel associated to \(((1+z)/(1-z))^d\).  This theorem gives
\[
        \Gamma_d\cong F_2,
        \qquad
        \PP\Gamma_d\cong F_2.
\]
Zariski density follows from the conjugates of the transvection, and infinite index follows from an elementary \(\Z^2\)-obstruction in the ambient arithmetic symplectic group.
\end{abstract}

\maketitle

\section{The normal form}

Let
\[
        d=2M\ge4,
        \qquad
        N=d-1=2M-1.
\]
Define \(c_d(n)\), \(n\ge0\), by
\[
        \sum_{n\ge0}c_d(n)z^n
        =
        \left(\frac{1+z}{1-z}\right)^d.
\]
Thus \(c_d(n)>0\) for \(n>0\), since
\[
        \left(\frac{1+z}{1-z}\right)^d
        =
        (1+z)^d(1-z)^{-d}
\]
is a product of a polynomial with nonnegative coefficients and a power series with strictly positive coefficients.

Define the odd kernel \(a_d:\Z\to\Z\) by
\[
        a_d(0)=0,
        \qquad
        a_d(n)=c_d(n)\quad(n>0),
        \qquad
        a_d(-n)=-c_d(n)\quad(n>0).
\]
For \(m\in\Z\), put
\[
        \rho_m(t)=\left(\frac{1+t}{1-t}\right)^m
\]
as a formal power series at \(t=0\).  Define the odd coefficient functional
\[
        \Lambda(F)=2\sum_{j=0}^{M-1}[t^{2j+1}]F(t).
\]

Throughout, \(\Comp(h)\) denotes the column companion matrix: if
\[
        h(x)=x^d+a_{d-1}x^{d-1}+\cdots+a_1x+a_0,
\]
then
\[
        \Comp(h)e_i=e_{i+1}\quad(0\le i<d-1),
        \qquad
        \Comp(h)e_{d-1}=-\sum_{i=0}^{d-1}a_i e_i.
\]

\begin{lemma}\label{lem:kernel-polynomial}
For every \(n\in\Z\),
\[
        a_d(n)=\Lambda(\rho_n).
\]
Consequently \(a_d(n)\) is the restriction to \(\Z\) of an odd polynomial \(A_d(n)\in\Q[n]\) of degree \(N=d-1\), with leading coefficient
\[
        \frac{2^d}{(d-1)!}.
\]
In particular,
\[
        \sum_{r=0}^{d}(-1)^r\binom dr\,a_d(n-r)=0
        \qquad(n\in\Z).
\]
\end{lemma}

\begin{proof}
We first prove \(a_d(n)=\Lambda(\rho_n)\) without contour language.  In the ring of formal power series in \(z\), with coefficients in \(\Q[[t]]\), one has
\[
        \sum_{n\ge1}\rho_n(t)z^n
        =
        \frac{z\rho_1(t)}{1-z\rho_1(t)}
        =
        \frac{z(1+t)}{(1-t)-z(1+t)}.
\]
For \(r\ge1\), the coefficient of \(t^r\) in this series is
\[
        \frac{2z(1+z)^{r-1}}{(1-z)^{r+1}}.
\]
Hence
\[
\begin{aligned}
        \sum_{n\ge1}\Lambda(\rho_n)z^n
        &=
        2\sum_{j=0}^{M-1}
        [t^{2j+1}]
        \frac{z(1+t)}{(1-t)-z(1+t)}                         \\
        &=
        4z\sum_{j=0}^{M-1}
        \frac{(1+z)^{2j}}{(1-z)^{2j+2}}                      \\
        &=
        \frac{4z}{(1-z)^2}
        \sum_{j=0}^{M-1}
        \left(\frac{1+z}{1-z}\right)^{2j}.
\end{aligned}
\]
Since
\[
        \left(\frac{1+z}{1-z}\right)^2-1
        =
        \frac{4z}{(1-z)^2},
\]
the last finite geometric sum telescopes to
\[
        \left(\frac{1+z}{1-z}\right)^{2M}-1
        =
        \left(\frac{1+z}{1-z}\right)^d-1
        =
        \sum_{n\ge1}c_d(n)z^n.
\]
Thus \(\Lambda(\rho_n)=c_d(n)\) for every \(n>0\).  For \(n=0\), both sides are zero.  For \(n<0\), use
\[
        \rho_{-n}(t)=\rho_n(-t),
\]
and the fact that \(\Lambda\) reads only odd coefficients.  Hence
\[
        \Lambda(\rho_{-n})=-\Lambda(\rho_n),
\]
which is exactly the defining oddness of \(a_d\).

For fixed \(0\le r\le N\), the coefficient \([t^r]\rho_n(t)\) is a polynomial in \(n\) of degree at most \(r\), since
\[
        \rho_n(t)
        =
        \exp\left(n\log\frac{1+t}{1-t}\right).
\]
Therefore \(a_d(n)=\Lambda(\rho_n)\) is polynomial in \(n\) of degree at most \(N\).  It is odd because \(a_d(-n)=-a_d(n)\).  Its leading term comes only from the coefficient of \(t^N\).  Since
\[
        \log\frac{1+t}{1-t}=2t+O(t^3),
\]
we have
\[
        [t^N]\rho_n(t)
        =
        \frac{2^N}{N!}n^N+O(n^{N-1}).
\]
The outer factor \(2\) in \(\Lambda\) gives leading coefficient
\[
        \frac{2^{N+1}}{N!}
        =
        \frac{2^d}{(d-1)!}.
\]
The finite-difference identity is the vanishing of the \(d\)-th finite difference of a polynomial of degree \(d-1\).
\end{proof}

Let \(V\) be a \(d\)-dimensional rational vector space with basis
\[
        e_0,e_1,\ldots,e_N.
\]
Define \(A\in\GL(V)\) by the companion rule for \((x-1)^d\):
\[
        Ae_i=e_{i+1}\quad(0\le i<N),
\]
and, since \(d\) is even,
\[
        Ae_N=\sum_{j=0}^{N}(-1)^{j+1}\binom dj e_j.
\]
Thus \(A\) is one regular unipotent Jordan block and
\[
        (A-I)^d=0.
\]
Define an alternating bilinear form \(\Omega=\Omega_d\) on \(V\) by
\[
        \Omega(e_i,e_j)=a_d(j-i).
\]

\begin{proposition}\label{prop:symplectic-normal}
The form \(\Omega\) is nondegenerate, and \(A\) preserves \(\Omega\).
\end{proposition}

\begin{proof}
The form is alternating because \(a_d\) is odd.  We first prove \(A\)-invariance.  If \(i,j<N\), then
\[
        \Omega(Ae_i,Ae_j)
        =
        \Omega(e_{i+1},e_{j+1})
        =
        a_d(j-i)
        =
        \Omega(e_i,e_j).
\]
If \(j<N\), Lemma~\ref{lem:kernel-polynomial}, applied with \(n=j+1\), gives
\[
        \sum_{r=0}^{d}(-1)^r\binom dr\,a_d(j+1-r)=0.
\]
The omitted \(r=d\) term in the partial sum \(0\le r\le N=d-1\) is
\[
        a_d(j+1-d)=a_d(j-N),
\]
because \(d\) is even.  Hence
\[
        \sum_{r=0}^{N}(-1)^r\binom dr\,a_d(j+1-r)
        =
        -a_d(j-N).
\]
Therefore
\[
\begin{aligned}
        \Omega(Ae_N,Ae_j)
        &=
        \sum_{r=0}^{N}(-1)^{r+1}\binom dr\,a_d(j+1-r) \\
        &=
        a_d(j-N)
        =
        \Omega(e_N,e_j).
\end{aligned}
\]
The case \(\Omega(Ae_i,Ae_N)=\Omega(e_i,e_N)\) follows by alternation, and
\[
        \Omega(Ae_N,Ae_N)=0=\Omega(e_N,e_N).
\]
Thus \(A\) preserves \(\Omega\).

It remains to prove nondegeneracy.  Let \(\mathcal R\) be the radical of \(\Omega\).  Since \(A\) preserves \(\Omega\), the radical is \(A\)-stable.  If \(\mathcal R\ne0\), then, because \(A\) is a single unipotent Jordan block, \(\mathcal R\) contains the fixed line of \(A\), namely the line spanned by
\[
        \eta=(A-I)^Ne_0.
\]
Indeed, if \(0\ne w\in\mathcal R\), choose \(m\) maximal such that \((A-I)^mw\ne0\); then \((A-I)^mw\) is fixed by \(A\), and the fixed space of a single Jordan block is one-dimensional.

But
\[
\begin{aligned}
        \Omega(e_0,\eta)
        &=
        \sum_{r=0}^{N}(-1)^{N-r}\binom Nr\,a_d(r)        \\
        &=
        N!\,[n^N]A_d(n)
        =
        2^d\ne0.
\end{aligned}
\]
Thus \(\eta\notin\mathcal R\), a contradiction.  Therefore \(\Omega\) is nondegenerate.
\end{proof}

Define the symplectic transvection
\[
        T(x)=x+\Omega(x,e_0)e_0.
\]
Since \(\Omega(e_0,e_0)=0\), a direct computation gives
\[
        \Omega(Tx,Ty)=\Omega(x,y)
        \qquad(x,y\in V).
\]
Thus \(T\in\Sp(\Omega)\).  Moreover \(T\) is integral in the basis \(e_0,\ldots,e_N\), because
\[
        T(e_j)=e_j+\Omega(e_j,e_0)e_0
        =
        e_j+a_d(-j)e_0.
\]

\begin{lemma}\label{lem:orbit-gram}
For all \(p,q\in\Z\),
\[
        \Omega(A^pe_0,A^qe_0)=a_d(q-p).
\]
\end{lemma}

\begin{proof}
By \(A\)-invariance, it is enough to prove
\[
        \Omega(e_0,A^ne_0)=a_d(n)
        \qquad(n\in\Z).
\]
For \(n\ge0\), set
\[
        s(n)=\Omega(e_0,A^ne_0).
\]
The relation \((A-I)^d=0\) implies that \(s(n)\) satisfies the forward finite-difference equation
\[
        \sum_{r=0}^{d}(-1)^{d-r}\binom dr\,s(n+r)=0.
\]
The polynomial sequence \(a_d(n)\) satisfies the same equation by Lemma~\ref{lem:kernel-polynomial}.  The two sequences agree for \(0\le n\le N\), by the definition of \(\Omega\).  Hence
\[
        s(n)=a_d(n)
        \qquad(n\ge0).
\]
For \(n=-m<0\), \(A\)-invariance and skew-symmetry give
\[
        \Omega(e_0,A^{-m}e_0)
        =
        \Omega(A^me_0,e_0)
        =
        -\Omega(e_0,A^me_0)
        =
        -a_d(m)
        =
        a_d(-m).
\]
The two-variable formula follows by applying \(A^{-p}\) to both arguments.
\end{proof}

\begin{lemma}[Rank-one Levelt uniqueness]\label{lem:levelt}
Let \(f,g\in\Q[x]\) be coprime monic polynomials of degree \(d\).  Let \((A,B)\) and \((A',B')\) be pairs of rational linear maps on \(d\)-dimensional rational vector spaces such that
\[
        \chi_A=\chi_{A'}=f,
        \qquad
        \chi_B=\chi_{B'}=g,
        \qquad
        \rank(B-A)=\rank(B'-A')=1.
\]
Then the two pairs are rationally conjugate.
\end{lemma}

\begin{proof}
It is enough to put one pair \((A,B)\) into a basis depending only on \(f\) and \(g\).  Choose
\[
        0\ne u\in\im(B-A),
\]
and write
\[
        B=A+u\otimes\varphi,
        \qquad
        (u\otimes\varphi)(x)=\varphi(x)u,
\]
for a rational covector \(\varphi\).  Let
\[
        W=\Span_{\Q}\{A^iu:i\ge0\}.
\]
Then \(W\) is \(A\)-invariant.  It is also \(B\)-invariant because
\[
        (B-A)V\subseteq \Q u\subseteq W.
\]
On \(V/W\), the induced maps of \(A\) and \(B\) are equal.  Therefore the characteristic polynomial of this common quotient map divides both \(f\) and \(g\).  Since \(f\) and \(g\) are coprime, the quotient is zero.  Hence
\[
        u,Au,\ldots,A^{d-1}u
\]
is a basis of \(V\).

In this cyclic basis, the matrix of \(A\) is the column companion matrix of \(f\).  The matrix of \(B\) is determined by the scalars
\[
        \alpha_i=\varphi(A^iu),
        \qquad 0\le i\le d-1.
\]
The matrix determinant lemma gives, as a Laurent expansion at \(\lambda=\infty\),
\[
\begin{aligned}
        \frac{g(\lambda)}{f(\lambda)}
        &=
        \frac{\det(\lambda I-B)}{\det(\lambda I-A)}       \\
        &=
        1-\varphi((\lambda I-A)^{-1}u)                    \\
        &=
        1-\sum_{i=0}^{d-1}\alpha_i\lambda^{-i-1}
        +O(\lambda^{-d-1}).
\end{aligned}
\]
Thus \(\alpha_0,\ldots,\alpha_{d-1}\) are determined by the first \(d\) Laurent coefficients of \(1-g/f\) at infinity.  They depend only on \(f\) and \(g\).  Therefore any two pairs satisfying the hypotheses are rationally conjugate.
\end{proof}

\begin{proposition}\label{prop:B-characteristic}
Let
\[
        B=TA.
\]
Then
\[
        \chi_A(x)=(x-1)^d,
        \qquad
        \chi_B(x)=(x+1)^d.
\]
Moreover, the pair \((A,B)\) is rationally conjugate to the companion pair
\[
        \left(\Comp((x-1)^d),\Comp((x+1)^d)\right).
\]
\end{proposition}

\begin{proof}
The assertion for \(A\) is part of the construction.  Since
\[
        Bx=Ax+\Omega(Ax,e_0)e_0,
\]
we have
\[
        B=A+e_0\otimes\psi,
        \qquad
        \psi(x)=\Omega(Ax,e_0).
\]
Thus \(B-A\) has rank one.  The matrix determinant lemma gives
\[
        \det(\lambda I-B)
        =
        \det(\lambda I-A)
        \left(
        1-\psi((\lambda I-A)^{-1}e_0)
        \right).
\]
As a Laurent expansion at infinity,
\[
\begin{aligned}
        \psi((\lambda I-A)^{-1}e_0)
        &=
        \sum_{m\ge0}\lambda^{-m-1}\Omega(A^{m+1}e_0,e_0) \\
        &=
        \sum_{n\ge1}a_d(-n)\lambda^{-n}                  \\
        &=
        -\sum_{n\ge1}c_d(n)\lambda^{-n},
\end{aligned}
\]
by Lemma~\ref{lem:orbit-gram}.  Therefore
\[
\begin{aligned}
        1-\psi((\lambda I-A)^{-1}e_0)
        &=
        1+\sum_{n\ge1}c_d(n)\lambda^{-n}                 \\
        &=
        \left(\frac{1+\lambda^{-1}}{1-\lambda^{-1}}\right)^d
        =
        \left(\frac{\lambda+1}{\lambda-1}\right)^d.
\end{aligned}
\]
Since
\[
        \det(\lambda I-A)=(\lambda-1)^d,
\]
we obtain
\[
        \det(\lambda I-B)=(\lambda+1)^d.
\]

The two column companion matrices of \((x-1)^d\) and \((x+1)^d\) differ only in their final column, and therefore their difference has rank one.  The polynomials \((x-1)^d\) and \((x+1)^d\) are coprime.  Lemma~\ref{lem:levelt} gives the claimed rational conjugacy.
\end{proof}

\section{The transvection certificate}

Let
\[
        \mathcal C
        =
        \left\{
        Q(t)=q_0+q_1t+\cdots+q_Nt^N\in\mathbb R[t]_{\le N}
        :
        |q_0|>\sum_{n=1}^{N}|q_n|
        \right\}.
\]

\begin{lemma}[Cone renewal]\label{lem:cone-renewal}
Let \(Q\in\mathcal C\), let \(k,\mu\in\Z\setminus\{0\}\), and set
\[
        P=[\rho_kQ]_{\le N}=\sum_{n=0}^{N}p_nt^n,
        \qquad
        F=\Lambda(P).
\]
Then
\[
        F\ne0
\]
and
\[
        P-\mu F\in\mathcal C,
\]
where \(P-\mu F\) means that the constant coefficient of \(P\) is changed by subtracting \(\mu F\).
\end{lemma}

\begin{proof}
The cone \(\mathcal C\) is invariant under multiplication by \(-1\).  Moreover, replacing \(Q\) by \(-Q\) replaces \(P\) and \(F\) by \(-P\) and \(-F\), so both conclusions are unchanged up to the same overall sign.  Hence we may assume
\[
        q_0>0,
        \qquad
        q_0>\sum_{n=1}^{N}|q_n|.
\]

Write
\[
        k=\tau a,
        \qquad
        \tau\in\{\pm1\},
        \qquad
        a\ge1,
\]
and
\[
        \rho_a(t)=\sum_{n\ge0}C_a(n)t^n.
\]
The coefficients \(C_a(n)\) are positive and nondecreasing.  Positivity follows from
\[
        \rho_a(t)=\frac{(1+t)^a}{(1-t)^a},
\]
and monotonicity follows because
\[
        (1-t)\rho_a(t)=\frac{(1+t)^a}{(1-t)^{a-1}}
\]
has nonnegative coefficients.

Since
\[
        \rho_k(t)=\rho_a(\tau t)
        =
        \sum_{n\ge0}\tau^nC_a(n)t^n,
\]
we have
\[
        p_n=\sum_{j=0}^{n}\tau^{n-j}C_a(n-j)q_j.
\]
Define
\[
        V_n=\tau^np_n.
\]
Then
\[
        V_n
        =
        C_a(n)q_0+
        \sum_{j=1}^{n}\tau^jC_a(n-j)q_j.
\]
Using \(C_a(n-j)\le C_a(n)\), we obtain
\[
        V_n
        \ge
        C_a(n)\left(q_0-\sum_{j=1}^{n}|q_j|\right)>0.
\]
Thus
\[
        |p_n|=V_n
        \qquad(0\le n\le N).
\]

Define
\[
        T_a(m)=(-1)^m\sum_{r=0}^{m}(-1)^rC_a(r)
        \qquad(m\ge0).
\]
Then \(T_a(m)>0\), and \(T_a(m)\) is nondecreasing.  Indeed,
\[
        \sum_{m\ge0}T_a(m)t^m
        =
        \frac{\rho_a(t)}{1+t}
        =
        \frac{(1+t)^{a-1}}{(1-t)^a},
\]
which has positive coefficients; multiplying by \(1-t\) gives
\[
        \frac{(1+t)^{a-1}}{(1-t)^{a-1}},
\]
which has nonnegative coefficients.

Let
\[
        O=\sum_{\substack{1\le n\le N\\ n\text{ odd}}}V_n,
        \qquad
        E=\sum_{\substack{0\le n\le N\\ n\text{ even}}}V_n.
\]
Because \(N\) is odd,
\[
\begin{aligned}
        E-O
        &=
        \sum_{n=0}^{N}(-1)^nV_n                                      \\
        &=
        -q_0T_a(N)
        -
        \sum_{j=1}^{N}\tau^jq_jT_a(N-j).
\end{aligned}
\]
Since \(0<T_a(N-j)\le T_a(N)\), we get
\[
        E-O
        \le
        -T_a(N)\left(q_0-\sum_{j=1}^{N}|q_j|\right)<0.
\]
Therefore
\[
        O>E.
\]

For odd \(n\), one has \(p_n=\tau V_n\).  Hence
\[
        F=\Lambda(P)
        =
        2\sum_{\substack{1\le n\le N\\ n\text{ odd}}}p_n
        =
        2\tau O.
\]
Thus \(F\ne0\), and
\[
        |F|=2O>O+E=\sum_{n=0}^{N}|p_n|.
\]
The polynomial \(P-\mu F\) has nonconstant coefficients \(p_1,\ldots,p_N\), while its constant coefficient is \(p_0-\mu F\).  Since \(|\mu|\ge1\),
\[
\begin{aligned}
        |p_0-\mu F|
        &\ge
        |\mu|\,|F|-|p_0|  \\
        &\ge
        |F|-|p_0|         \\
        &>
        \sum_{n=1}^{N}|p_n|.
\end{aligned}
\]
Thus \(P-\mu F\in\mathcal C\).
\end{proof}

\begin{theorem}[Transvection renewal nonvanishing]\label{thm:renewal}
Let
\[
        k_1,\ldots,k_\ell\in\Z\setminus\{0\},
        \qquad
        p_0=0,
        \qquad
        p_s=k_1+\cdots+k_s.
\]
Let
\[
        \mu_1,\ldots,\mu_{\ell-1}\in\Z\setminus\{0\}.
\]
Define \(D_s\), \(1\le s\le\ell\), recursively by
\[
        D_s
        =
        a_d(p_s)
        -
        \sum_{i=1}^{s-1}\mu_iD_i\,a_d(p_s-p_i),
\]
with the empty sum interpreted as zero.  Then
\[
        D_s\ne0
        \qquad(1\le s\le\ell).
\]
\end{theorem}

\begin{proof}
Put
\[
        Q_0=1.
\]
For \(1\le s\le\ell\), define
\[
        D_s=\Lambda(\rho_{k_s}Q_{s-1}),
\]
and, for \(s<\ell\), define
\[
        Q_s=[\rho_{k_s}Q_{s-1}-\mu_sD_s]_{\le N}.
\]
Since \(Q_0=1\in\mathcal C\), Lemma~\ref{lem:cone-renewal} implies inductively that
\[
        Q_s\in\mathcal C
        \quad(0\le s\le\ell-1),
        \qquad
        D_s\ne0
        \quad(1\le s\le\ell).
\]

It remains only to identify this recursion with the displayed formula.  We claim that, for \(0\le s\le\ell-1\),
\[
        Q_s
        =
        \left[
        \rho_{p_s}
        -
        \sum_{i=1}^{s}\mu_iD_i\rho_{p_s-p_i}
        \right]_{\le N}.
\]
This is immediate for \(s=0\).  If it holds for \(s-1\), then multiplication by \(\rho_{k_s}\), followed by subtracting \(\mu_sD_s\), gives exactly the formula for \(Q_s\).

Using Lemma~\ref{lem:kernel-polynomial}, and the fact that \(\Lambda\) only reads coefficients of degrees at most \(N\), we get
\[
\begin{aligned}
        D_s
        &=
        \Lambda(\rho_{k_s}Q_{s-1})                                      \\
        &=
        \Lambda\left(
        \rho_{p_s}
        -
        \sum_{i=1}^{s-1}\mu_iD_i\rho_{p_s-p_i}
        \right)                                                         \\
        &=
        a_d(p_s)
        -
        \sum_{i=1}^{s-1}\mu_iD_i\,a_d(p_s-p_i).
\end{aligned}
\]
Thus the recursively defined numbers \(D_s\) are exactly the displayed ones, and each is nonzero.
\end{proof}

\section{Freeness}

For \(v\in V\), write
\[
        \tau_v^m(x)=x+m\Omega(x,v)v
        \qquad(m\in\Z).
\]
Thus \(T^m=\tau_{e_0}^m\), and, because \(A\) preserves \(\Omega\),
\[
        A^pT^mA^{-p}=\tau_{A^pe_0}^m.
\]

\begin{lemma}[Word pairing]\label{lem:word-pairing}
Let
\[
        W=T^{m_1}A^{k_1}T^{m_2}A^{k_2}\cdots T^{m_\ell}A^{k_\ell},
        \qquad
        m_i,k_i\in\Z\setminus\{0\}.
\]
Then
\[
        \Omega(e_0,We_0)\ne0.
\]
\end{lemma}

\begin{proof}
Put
\[
        p_0=0,
        \qquad
        p_s=k_1+\cdots+k_s,
        \qquad
        v_s=A^{p_s}e_0.
\]
Then
\[
        W
        =
        \tau_{v_0}^{m_1}
        \tau_{v_1}^{m_2}
        \cdots
        \tau_{v_{\ell-1}}^{m_\ell}
        A^{p_\ell}.
\]
Therefore
\[
        We_0
        =
        \tau_{v_0}^{m_1}
        \tau_{v_1}^{m_2}
        \cdots
        \tau_{v_{\ell-1}}^{m_\ell}
        v_\ell.
\]
Since \(\tau_{v_0}^{m_1}\) fixes \(v_0=e_0\) and preserves \(\Omega\),
\[
        \Omega(e_0,We_0)
        =
        \Omega\left(v_0,
        \tau_{v_1}^{m_2}\cdots\tau_{v_{\ell-1}}^{m_\ell}v_\ell
        \right).
\]

Define
\[
        y_1=v_0,
        \qquad
        y_{s+1}=\tau_{v_s}^{-m_{s+1}}y_s
        \quad(1\le s\le\ell-1),
\]
and
\[
        D_s=\Omega(y_s,v_s)
        \qquad(1\le s\le\ell).
\]
Repeatedly moving transvections from the second argument to the first gives
\[
        \Omega(e_0,We_0)=D_\ell.
\]
Moreover, by induction,
\[
        y_s
        =
        v_0-\sum_{i=1}^{s-1}m_{i+1}D_i v_i.
\]
Indeed, the induction step is
\[
        y_{s+1}
        =
        y_s-m_{s+1}\Omega(y_s,v_s)v_s
        =
        y_s-m_{s+1}D_sv_s.
\]
Thus Lemma~\ref{lem:orbit-gram} gives
\[
\begin{aligned}
        D_s
        &=
        \Omega(y_s,v_s)                                               \\
        &=
        a_d(p_s)
        -
        \sum_{i=1}^{s-1}m_{i+1}D_i\,a_d(p_s-p_i).
\end{aligned}
\]
This is precisely the recurrence of Theorem~\ref{thm:renewal}, with
\[
        \mu_i=m_{i+1}.
\]
Hence \(D_\ell\ne0\).  Therefore
\[
        \Omega(e_0,We_0)\ne0.
\]
\end{proof}

\begin{proposition}\label{prop:free}
For even \(d\ge4\),
\[
        G_d:=\langle A,T\rangle\cong F_2,
        \qquad
        \PP G_d\cong F_2.
\]
\end{proposition}

\begin{proof}
Let \(F_2=\langle a,t\rangle\), and map
\[
        a\mapsto A,
        \qquad
        t\mapsto T.
\]
It is enough to prove that no nontrivial reduced word maps to a scalar matrix.

Pure powers are not projectively scalar.  If \(k\ne0\), then \(A^k\) is a nontrivial unipotent because \(A\) is one regular unipotent Jordan block.  Hence \(A^k\) cannot be scalar.  Similarly,
\[
        T^m=I+mN,
        \qquad
        N(x)=\Omega(x,e_0)e_0,
\]
is a nontrivial unipotent for \(m\ne0\), and is not scalar.

Every nontrivial mixed reduced word is conjugate to a cyclically reduced word of the form
\[
        t^{m_1}a^{k_1}t^{m_2}a^{k_2}\cdots t^{m_\ell}a^{k_\ell},
        \qquad
        m_i,k_i\ne0.
\]
Projective scalarity is invariant under conjugation.  Let \(W\) be the image of such a word.  If \(W=\lambda I\), then
\[
        \Omega(e_0,We_0)=\lambda\Omega(e_0,e_0)=0,
\]
because \(e_0\) is isotropic.  This contradicts Lemma~\ref{lem:word-pairing}.  Therefore no nontrivial reduced word maps to a scalar matrix.

It follows that the homomorphism \(F_2\to G_d\) is injective, and that the induced homomorphism \(F_2\to\PP G_d\) is also injective.  Both maps are surjective by definition.  Hence
\[
        G_d\cong F_2,
        \qquad
        \PP G_d\cong F_2.
\]
\end{proof}

\section{Density and infinite index}

In the basis \(e_0,\ldots,e_N\), the matrices of \(A\), \(T\), and \(\Omega\) are integral.  Hence
\[
        G_d\le \Sp(\Omega,\Z),
\]
where
\[
        \Sp(\Omega,\Z)
        =
        \{g\in\GL_d(\Z):\Omega(gx,gy)=\Omega(x,y)\text{ for all }x,y\}.
\]

\begin{proposition}\label{prop:zariski}
The group \(G_d\) is Zariski dense in \(\Sp(\Omega)\).
\end{proposition}

\begin{proof}
We work after extension of scalars to \(\C\).  Let \(H\) be the Zariski closure of \(G_d\) in \(\Sp(\Omega)_\C\).  For \(0\le i\le N\), put
\[
        v_i=A^ie_0.
\]
These vectors form a basis of \(V_\C\).  The group \(G_d\) contains the conjugate transvections
\[
        T_i=A^iTA^{-i},
        \qquad
        T_i(x)=x+\Omega(x,v_i)v_i.
\]
Write
\[
        T_i=I+N_i,
        \qquad
        N_i(x)=\Omega(x,v_i)v_i.
\]
Since \(N_i^2=0\),
\[
        T_i^n=I+nN_i
        \qquad(n\in\Z).
\]
The powers of \(T_i\) are Zariski dense in the one-parameter unipotent subgroup
\[
        \{I+sN_i:s\in\C\}.
\]
Therefore
\[
        N_i\in\Lie(H)
        \qquad(0\le i\le N).
\]

For \(i\ne j\),
\[
        \Omega(v_i,v_j)=a_d(j-i)\ne0,
\]
because \(a_d(n)=\pm c_d(|n|)\) for \(n\ne0\), and \(c_d(n)>0\) for \(n>0\).  A direct computation gives
\[
        [N_i,N_j](x)
        =
        -\Omega(v_i,v_j)
        \left(
        \Omega(x,v_j)v_i+\Omega(x,v_i)v_j
        \right).
\]
Thus \(\Lie(H)\) contains the diagonal operators
\[
        x\mapsto \Omega(x,v_i)v_i
\]
and all off-diagonal symmetric operators
\[
        x\mapsto \Omega(x,v_j)v_i+\Omega(x,v_i)v_j.
\]
Under the standard identification
\[
        \operatorname{Sym}^2(V_\C)\cong\mathfrak{sp}(V_\C,\Omega),
        \qquad
        u\odot v\mapsto
        \left(x\mapsto \Omega(x,u)v+\Omega(x,v)u\right),
\]
these operators are the basis elements corresponding to the symmetric tensors in the basis
\[
        v_0,\ldots,v_N.
\]
Hence
\[
        \Lie(H)=\mathfrak{sp}(V_\C,\Omega).
\]
Since \(\Sp(\Omega)_\C\) is connected, this implies
\[
        H=\Sp(\Omega)_\C.
\]
Thus \(G_d\) is Zariski dense in \(\Sp(\Omega)\).
\end{proof}

\begin{lemma}\label{lem:ambient-Z2}
Every finite-index subgroup of \(\Sp(\Omega,\Z)\) contains a subgroup isomorphic to \(\Z^2\).
\end{lemma}

\begin{proof}
Let
\[
        X=\log A
        =
        \sum_{r=1}^{d-1}(-1)^{r+1}\frac{(A-I)^r}{r}.
\]
This is a rational nilpotent matrix, and \(A=\exp X\).  For every integer \(n\),
\[
        A^n=\exp(nX)
\]
preserves \(\Omega\).  The entries of
\[
        \exp(tX)^{\mathsf t}\Omega\exp(tX)-\Omega
\]
are polynomials in \(t\), since \(X\) is nilpotent.  They vanish for all \(t\in\Z\), hence vanish identically.  Differentiating at \(t=0\) gives
\[
        X\in\mathfrak{sp}(\Omega)(\Q).
\]
If \(Y\in\mathfrak{sp}(\Omega)\), then every odd power of \(Y\) again lies in \(\mathfrak{sp}(\Omega)\), because
\[
        \Omega(Y^m x,y)=(-1)^m\Omega(x,Y^m y).
\]
Thus
\[
        X^3\in\mathfrak{sp}(\Omega)(\Q).
\]

Choose \(L\ge1\) sufficiently divisible that
\[
        E=\exp(LX^3)
\]
has integral entries.  This is possible because \(X^3\) is rational and nilpotent, so the exponential is a finite sum with rational coefficients.  Since \(LX^3\in\mathfrak{sp}(\Omega)\), the matrix \(E\) preserves \(\Omega\).  It is unipotent and integral, hence
\[
        E\in\Sp(\Omega,\Z).
\]
The elements
\[
        A=\exp X,
        \qquad
        E=\exp(LX^3)
\]
commute.

They generate a copy of \(\Z^2\).  Indeed, if
\[
        A^mE^n=I,
\]
then, since \(X\) and \(X^3\) commute,
\[
        \exp(mX+nLX^3)=I.
\]
The exponent \(mX+nLX^3\) is nilpotent, so applying the finite polynomial logarithm gives
\[
        mX+nLX^3=0.
\]
Since \(A\) is one Jordan block and \(d\ge4\), \(X^3\ne0\).  Moreover \(X\) and \(X^3\) are linearly independent: if \(N_A=A-I\), then
\[
        X=N_A+\text{higher powers of }N_A,
        \qquad
        X^3=N_A^3+\text{higher powers of }N_A.
\]
Therefore \(m=n=0\), and
\[
        \langle A,E\rangle\cong\Z^2.
\]

Now let \(K\le\Sp(\Omega,\Z)\) have finite index.  Some positive power of \(A\) lies in \(K\), and some positive power of \(E\) lies in \(K\).  These two positive powers still generate a copy of \(\Z^2\).  Hence every finite-index subgroup of \(\Sp(\Omega,\Z)\) contains a subgroup isomorphic to \(\Z^2\).
\end{proof}

\begin{lemma}\label{lem:free-no-Z2}
A free group contains no subgroup isomorphic to \(\Z^2\).
\end{lemma}

\begin{proof}
Let \(F\) be a free group with its Cayley tree.  A nontrivial element of \(F\) acts on this tree by translation along a unique axis.  If another element commutes with it, then it preserves this axis and acts on it by a translation.  The action of \(F\) on its Cayley tree is free, so the induced homomorphism from the centralizer to the infinite cyclic translation group of the axis is injective.  Hence the centralizer of every nontrivial element is cyclic.  Therefore every abelian subgroup of \(F\) is cyclic, and \(F\) contains no \(\Z^2\).
\end{proof}

\begin{proposition}\label{prop:normal-infinite-index}
The group \(G_d\) has infinite index in \(\Sp(\Omega,\Z)\).
\end{proposition}

\begin{proof}
By Proposition~\ref{prop:free},
\[
        G_d\cong F_2.
\]
Thus \(G_d\) contains no subgroup isomorphic to \(\Z^2\), by Lemma~\ref{lem:free-no-Z2}.  If \(G_d\) had finite index in \(\Sp(\Omega,\Z)\), then Lemma~\ref{lem:ambient-Z2} would force \(G_d\) to contain a copy of \(\Z^2\), contradiction.
\end{proof}

\section{Companion coordinates}

Let
\[
        C_+=\Comp((x-1)^d),
        \qquad
        C_-=\Comp((x+1)^d).
\]
By Proposition~\ref{prop:B-characteristic}, there exists \(P\in\GL_d(\Q)\) such that
\[
        PAP^{-1}=C_+,
        \qquad
        PBP^{-1}=C_-.
\]
Since \(B=TA\), we have
\[
        \langle A,B\rangle=\langle A,T\rangle=G_d.
\]
Therefore
\[
        \Gamma_d:=
        \left\langle
        \Comp((x-1)^d),\Comp((x+1)^d)
        \right\rangle
        =
        PG_dP^{-1}.
\]

Define the transported alternating form
\[
        \Omega_d^{\rm orig}=P^{-\mathsf t}\Omega P^{-1}.
\]
After multiplying by a nonzero integer, assume that \(\Omega_d^{\rm orig}\) has integral entries.  Multiplying an alternating form by a nonzero scalar does not change its symplectic group.  Since \(\Gamma_d\) preserves \(\Omega_d^{\rm orig}\) and its generators are integral, we have
\[
        \Gamma_d\le\Sp(\Omega_d^{\rm orig},\Z).
\]

We now transfer infinite index from the normal-form lattice to the companion lattice.  Let
\[
        \Lambda_0=\Z^d
\]
be the normal-form lattice with basis \(e_0,\ldots,e_N\), and let
\[
        \Lambda_1=P^{-1}\Z^d.
\]
For a lattice \(\Lambda\subset V\), write
\[
        \Sp(\Omega,\Lambda)
        =
        \{g\in\Sp(\Omega)(\Q):g\Lambda=\Lambda\}.
\]
Then
\[
        \Sp(\Omega,\Lambda_0)=\Sp(\Omega,\Z),
\]
and
\[
        \Sp(\Omega,\Lambda_1)
        =
        P^{-1}\Sp(\Omega_d^{\rm orig},\Z)P.
\]
The lattices \(\Lambda_0\) and \(\Lambda_1\) are commensurable.  Choose \(r\ge1\) such that
\[
        r\Lambda_0\subseteq\Lambda_1\subseteq r^{-1}\Lambda_0.
\]
There are only finitely many lattices \(L\) satisfying
\[
        r\Lambda_0\subseteq L\subseteq r^{-1}\Lambda_0.
\]
The group \(\Sp(\Omega,\Lambda_0)\) acts on this finite set, and the stabilizer of \(\Lambda_1\) is
\[
        \Sp(\Omega,\Lambda_0)\cap\Sp(\Omega,\Lambda_1).
\]
Hence this intersection has finite index in \(\Sp(\Omega,\Lambda_0)\).  Interchanging \(\Lambda_0\) and \(\Lambda_1\) gives finite index in \(\Sp(\Omega,\Lambda_1)\) as well.

The group \(G_d\) lies in both stabilizers.  It lies in \(\Sp(\Omega,\Lambda_0)\) by construction, and it lies in \(\Sp(\Omega,\Lambda_1)\) because
\[
        PG_dP^{-1}=\Gamma_d\le\GL_d(\Z).
\]
Put
\[
        K=
        \Sp(\Omega,\Lambda_0)\cap\Sp(\Omega,\Lambda_1).
\]
If \(\Gamma_d\) had finite index in \(\Sp(\Omega_d^{\rm orig},\Z)\), then conjugating by \(P^{-1}\) would make \(G_d\) finite index in \(\Sp(\Omega,\Lambda_1)\).  Since
\[
        G_d\subseteq K\subseteq\Sp(\Omega,\Lambda_1),
\]
we would have
\[
        [K:G_d]\le[\Sp(\Omega,\Lambda_1):G_d]<\infty.
\]
Because \(K\) has finite index in \(\Sp(\Omega,\Lambda_0)=\Sp(\Omega,\Z)\), this would imply that \(G_d\) has finite index in \(\Sp(\Omega,\Z)\), contradicting Proposition~\ref{prop:normal-infinite-index}.  Therefore
\[
        [\Sp(\Omega_d^{\rm orig},\Z):\Gamma_d]=\infty.
\]

Zariski density and freeness transfer immediately under rational conjugacy.  Combining the preceding results gives the main theorem.

\begin{theorem}\label{thm:main}
For every even \(d\ge4\), the even all-half hypergeometric group
\[
        \Gamma_d=
        \left\langle
        \Comp((x-1)^d),\Comp((x+1)^d)
        \right\rangle
\]
is thin in its ambient integral symplectic group.  More precisely, for the integral rational alternating form \(\Omega_d^{\rm orig}\) constructed above,
\[
        \overline{\Gamma_d}^{\,\mathrm{Zar}}
        =
        \Sp(\Omega_d^{\rm orig})
\]
and
\[
        [\Sp(\Omega_d^{\rm orig},\Z):\Gamma_d]=\infty.
\]
Moreover,
\[
        \Gamma_d\cong F_2,
        \qquad
        \PP\Gamma_d\cong F_2.
\]
\end{theorem}

\begin{proof}
The rational conjugacy
\[
        \Gamma_d=PG_dP^{-1}
\]
transfers Proposition~\ref{prop:free} to
\[
        \Gamma_d\cong F_2,
        \qquad
        \PP\Gamma_d\cong F_2.
\]
It transfers Proposition~\ref{prop:zariski} to
\[
        \overline{\Gamma_d}^{\,\mathrm{Zar}}
        =
        \Sp(\Omega_d^{\rm orig}).
\]
The infinite-index statement was proved in the commensurability argument above.  Thus \(\Gamma_d\) is Zariski dense and of infinite index in its ambient integral symplectic group.
\end{proof}

\begin{remark}
The case \(d=2\) is not included.  The proof of infinite index uses the independent commuting unipotents \(\exp(X)\) and \(\exp(LX^3)\), where \(X=\log A\); this requires \(X^3\ne0\), hence \(d\ge4\).  The even-dimensional tower is symplectic because the kernel \(a_d\) is odd and the rank-one local monodromy \(BA^{-1}\) is a transvection, whereas the odd-dimensional all-half tower is orthogonal and the corresponding local monodromy is a reflection.
\end{remark}

\begin{remark}
The theorem gives a uniform proof across the even all-half tower.  In particular, the same certificate specializes to the previously known small even cases, including the \(d=4\) Brav--Thomas case and the \(d=6\) Bajpai--Dona--Nitsche \(A\)-1 case.
\end{remark}

\end{document}